\chapter{Galerkin Discretisation and Isoparametric Elements}
\label{ch:discretization}

We now turn the weak form~\eqref{eq:weak} into a finite system of
algebraic equations. The idea is Galerkin's: replace the
infinite-dimensional spaces $\mathcal S,\mathcal V$ with
finite-dimensional subspaces built from piecewise polynomials over a
mesh of \emph{elements}. This chapter constructs the element
interpolation, the strain--displacement operator, and the isoparametric
mapping that makes curved, distorted elements practical.

\section{Interpolation and shape functions}

Partition $\domain$ into non-overlapping elements $\domain^e$. Over each
element the displacement field is interpolated from \emph{nodal} values
$\dofe$ using \textbf{shape (basis) functions} $N_a$:
\begin{equation}
  \disp^h(\vec x) = \sum_{a=1}^{n_{en}} N_a(\vec x)\,\vec d_a
                  = \Nmat(\vec x)\,\dofe ,
  \label{eq:interp}
\end{equation}
where $n_{en}$ is the number of element nodes and $\Nmat$ arranges the
$N_a$ into the shape-function matrix. Shape functions satisfy the
\textbf{Kronecker-delta / partition-of-unity} properties
\begin{equation}
  N_a(\vec x_b)=\delta_{ab}, \qquad \sum_a N_a(\vec x)=1,
\end{equation}
so nodal coefficients \emph{are} nodal displacements and rigid-body
translations are representable exactly (a prerequisite for convergence,
\cref{ch:meshing-convergence}).

\section{The strain--displacement matrix \texorpdfstring{$\Bmat$}{B}}

Applying the symmetric gradient~\eqref{eq:strain-disp} to the
interpolation~\eqref{eq:interp} gives the strain in terms of nodal DOF:
\begin{equation}
  \evec = \grad_{\!s}\disp^h = \Bmat\,\dofe, \qquad \Bmat=\grad_{\!s}\Nmat,
  \label{eq:Bdef}
\end{equation}
so $\Bmat$ contains shape-function derivatives. In 2D the block for node
$a$ is
\begin{equation}
  \Bmat_a =
  \begin{bmatrix}
    \partial N_a/\partial x & 0 \\
    0 & \partial N_a/\partial y \\
    \partial N_a/\partial y & \partial N_a/\partial x
  \end{bmatrix},
\end{equation}
with rows corresponding to $(\varepsilon_{xx},\varepsilon_{yy},\gamma_{xy})$;
the 3D case is the direct six-row analogue. $\Bmat$ is the single most
important operator in the FEM: it maps discrete DOF to strains and,
through $\ke=\int\Bmat\transpose\Dmat\,\Bmat\dd V$
(\cref{ch:elements-assembly}), to element stiffness.

\section{Isoparametric mapping}

Rather than compute $N_a(\vec x)$ directly for every distorted element,
we define the shape functions once on a fixed \textbf{parent element} in
\emph{natural coordinates} $\vec\xi=(\xi,\eta,\zeta)$ and map to physical
space. The \textbf{isoparametric} principle uses the \emph{same} shape
functions for geometry and field:
\begin{equation}
  \vec x = \sum_a N_a(\vec\xi)\,\vec x_a, \qquad
  \disp^h = \sum_a N_a(\vec\xi)\,\vec d_a .
  \label{eq:isoparam}
\end{equation}
Derivatives transform through the \textbf{Jacobian} of the map,
\begin{equation}
  \Jmat = \pd{\vec x}{\vec\xi} = \sum_a \pd{N_a}{\vec\xi}\,\vec x_a\transpose,
  \qquad
  \pd{N_a}{\vec x} = \Jmat^{-1}\pd{N_a}{\vec\xi},
  \qquad
  \dd V = \det\Jmat \,\dd\xi\,\dd\eta\,\dd\zeta .
  \label{eq:jacobian}
\end{equation}
A valid map requires $\det\Jmat>0$ at every integration point; a
non-positive Jacobian signals an inverted or excessively distorted
element and is a common cause of solver failure.

\section{A catalogue of common elements}

Shape functions for the workhorse elements (natural coordinates
$\xi,\eta,\zeta\in[-1,1]$ for quad/hex families; area/volume
\emph{barycentric} coordinates $L_i$ for triangle/tet families):

\begin{center}
\begin{tabularx}{\linewidth}{p{3.1cm}p{1.1cm}X}
  \toprule
  Element & Nodes & Shape functions \\
  \midrule
  Bar (linear)     & 2 & $N_{1,2}=\tfrac12(1\mp\xi)$ \\
  Bar (quadratic)  & 3 & $N_1=\tfrac12\xi(\xi-1),\ N_2=1-\xi^2,\ N_3=\tfrac12\xi(\xi+1)$ \\
  Tri (linear, CST)& 3 & $N_i=L_i$ (constant strain) \\
  Tri (quadratic)  & 6 & corner $L_i(2L_i-1)$, midside $4L_iL_j$ \\
  Quad (bilinear)  & 4 & $N_a=\tfrac14(1+\xi\xi_a)(1+\eta\eta_a)$ \\
  Quad (serendip.) & 8 & corner $\tfrac14(1+\xi\xi_a)(1+\eta\eta_a)(\xi\xi_a+\eta\eta_a-1)$ \\
  Tet (linear)     & 4 & $N_i=L_i$ (constant strain) \\
  Tet (quadratic)  & 10& corner $L_i(2L_i-1)$, midedge $4L_iL_j$ \\
  Hex (trilinear)  & 8 & $N_a=\tfrac18(1+\xi\xi_a)(1+\eta\eta_a)(1+\zeta\zeta_a)$ \\
  Hex (serendip.)  & 20& corner \& midedge serendipity formulae \\
  \bottomrule
\end{tabularx}
\end{center}

\begin{remark}[Linear vs quadratic, and why FreeCAD defaults to Tet10]
Linear simplex elements (Tri3, Tet4) have \emph{constant} strain over the
element: they capture bending only through many elements and are prone to
shear locking. Quadratic elements (Tri6, Tet10) represent linear strain
fields, model curved boundaries far better, and are dramatically more
accurate per node for stress analysis. This is why the FreeCAD/CalculiX
default for solids is the 10-node tetrahedron (C3D10), discussed in
\cref{ch:solvers}. Automatic tetrahedral meshers fill arbitrary solids
that hexahedra cannot, at the cost of more elements for equal accuracy.
\end{remark}

\begin{remark}[Serendipity vs Lagrange]
The 8-node quad and 20-node hex are \emph{serendipity} elements: only
edge nodes, incomplete high-order polynomial basis. Their Lagrange
cousins (Quad9, Hex27) add interior nodes for a complete basis and
slightly better accuracy, at higher cost. Either family requires
matching the correct quadrature rule (\cref{ch:elements-assembly}).
\end{remark}

With the interpolation, the $\Bmat$ operator, and the isoparametric map
in hand, the next chapter integrates $\Bmat\transpose\Dmat\Bmat$ to form
element stiffness, assembles the global system, and solves it.
